\appendix

\section{Extended Experimental Details}\label{app:experiments}

\subsection{Initialization Ablation}

We test whether the diagonal-dominance fingerprint depends on initialization scheme. We train depth-24 residual MLPs ($h{=}64$, $d{=}24$) on a synthetic regression target for 200 epochs from seven initialization schemes: Kaiming uniform/normal, Xavier uniform/normal, Gaussian, uniform, and orthogonal. Table~\ref{tab:init-ablation-full} reports pair accuracy before and after training.

\begin{table}[h]
\centering
\small
\begin{tabular}{@{}lccc@{}}
\toprule
Init scheme & Untrained & Trained & AUC \\
\midrule
Orthogonal       & $0.0\%$  & $100\%$ & $0.947$ \\
Kaiming-normal   & $2.1\%$  & $97\%$  & $0.981$ \\
Kaiming-uniform  & $2.1\%$  & $98.6\%$ & $0.98$ \\
Xavier-normal    & $2.1\%$  & $100\%$ & $0.990$ \\
Xavier-uniform   & $2.1\%$  & $100\%$ & $0.99$ \\
Uniform          & $2.1\%$  & $93.8\%$ & --- \\
Gaussian $\sigma{=}0.02$ & $2.1\%$ & $13\%$ & $0.671$ \\
\bottomrule
\end{tabular}
\caption{Initialization ablation on depth-24 residual MLPs (3 seeds, 200 epochs). All schemes show chance-level accuracy (${\leq}2.1\%$) at initialization, rising to $75$--$100\%$ after training. Orthogonal initialization---which provides dynamical isometry at init---shows \emph{zero} correct pairs before training.}
\label{tab:init-ablation-full}
\end{table}

Critically, orthogonal initialization provides dynamical isometry at initialization by construction~\citep{saxe2014exact}, yet shows \emph{zero} correct pairs before training, ruling out the hypothesis that near-orthogonal Jacobians alone create the fingerprint. The $\sigma{=}0.02$ small-Gaussian case fails because residual blocks collapse to near-zero contribution; without active per-block contribution there is no Jacobian constraint to enforce.

\subsection{PlainNet Control Experiment}

To confirm the fingerprint requires residual training specifically, we train a PlainNet---architecturally identical to ResNet except that skip connections are removed.

\begin{table}[h]
\centering
\small
\begin{tabular}{@{}lcccc@{}}
\toprule
Architecture & Eval Loss & Pair Acc & AUC & Neg Trace \\
\midrule
ResNet-24   & 1.35 & $100\%$ & $0.98$ & $100\%$ \\
PlainNet-24 & 1.56 & $3\%$   & $0.51$ & $47\%$ \\
\bottomrule
\end{tabular}
\caption{ResNet vs.\ PlainNet control. Despite similar evaluation loss, PlainNet achieves only chance-level pair accuracy (3\% vs.\ chance $\approx 4.2\%$), confirming the skip connection is necessary for the fingerprint.}
\label{tab:plainnet-control}
\end{table}

While both networks converge to similar evaluation loss (1.35 vs 1.56), ResNet achieves 100\% pair accuracy with 100\% of correctly paired products exhibiting negative trace, while PlainNet achieves only 3\%---at the chance baseline---with AUC 0.51 and no discernible diagonal structure. The skip connection creates a functional constraint that imprints the characteristic trace structure.

\subsection{Training Dynamics}

The fingerprint emerges rapidly during training. On a 48-block residual network:
\begin{itemize}[nosep]
\item Epoch 0: Mean correct-pair score = 0.19, pair accuracy = 6.25\%
\item Epoch 5: Diagonal fully separated, pair accuracy = 100\%
\item Epoch 300: Mean correct-pair score = 2.07, mean incorrect-pair score = 0.16 ($21\times$ ratio)
\end{itemize}
The transformation occurs within the first few epochs, well before training loss plateaus.

\section{Extended Theoretical Results}\label{app:proofs}

\subsection{Null Model (Corollary)}

\begin{corollary}[Null Model]
For randomly initialized weights with no training, all pairs satisfy $\mathbb{E}[s(i,j)] = 1/\sqrt{d} + o(1)$ uniformly in $(i,j)$. Hungarian matching recovers correct pairs at the chance rate $1/L$.
\end{corollary}

The null model rules out three alternative explanations:
\begin{enumerate}[nosep]
\item Diagonal dominance is not an artifact of compatible matrix shapes
\item The residual connection itself does not create the signal without training
\item Any nontrivial loss level does not suffice---training must actively couple the projections
\end{enumerate}

The diagonal-dominance fingerprint is a \emph{learned} property induced by gradient descent under the residual constraint.

\subsection{Geometric Interpretation of the Score}

The diagonal-dominance score $s(M) = |\mathrm{tr}(M)|/\|M\|_F$ has a natural geometric interpretation:
\begin{itemize}[nosep]
\item A matrix with energy uniformly spread across entries achieves $s = \Theta(1/\sqrt{d})$
\item A matrix proportional to the identity achieves $s = \sqrt{d}$
\end{itemize}

The key insight is that correctly paired blocks---where all $K$ matrices come from the same residual branch---score at $\Theta(\sqrt{d})$, while incorrect pairings remain at the random baseline $\Theta(1/\sqrt{d})$. This yields a signal-to-noise gap that grows linearly with hidden dimension $d$, making correct pairs increasingly separable in wider networks.

For typical trained blocks, $\|E\|_F/(\varepsilon\sqrt{d}) \in [1.9, 3.8]$, yielding $s(i,i) \in [1.76, 3.23]$---well above the baseline but below the theoretical maximum of $\sqrt{d}$.

\section{Verification Protocol Details}\label{app:verification}

\subsection{Gated Branch Score}

Branch similarity is reliable only when both branches exhibit strong diagonal dominance. We gate the cosine similarity:
\begin{equation}
G(\phi^A_\ell, \phi^B_m) = \langle \phi^A_\ell, \phi^B_m \rangle \cdot \min\!\left(\tfrac{s(M^A_\ell)}{\tau_s}, \tfrac{s(M^B_m)}{\tau_s}, 1\right)
\end{equation}
where $\tau_s$ is the minimum diagonal-dominance score across reference branches.

\subsection{Statistical Calibration}

We calibrate against a null distribution $\mathcal{N}_A = \{\mathcal{L}(A, U_j) : U_j \notin \mathcal{D}(A)\}$ of scores between reference $A$ and independently trained models. The z-score is:
\begin{equation}
Z(A,B) = \frac{\mathcal{L}(A,B) - \mu_{\mathrm{null}}}{\sigma_{\mathrm{null}}}
\end{equation}
Under Gaussian approximation, $\tau = 1.645$ yields 95\% confidence; $\tau = 3.0$ yields FPR $< 0.13\%$ for high-stakes applications.

\subsection{Full Protocol Steps}

\begin{enumerate}[nosep,leftmargin=*]
\item \textbf{Compatibility Check:} Verify $\mathrm{arch}(A) = \mathrm{arch}(B)$. If not, return \textsc{Incompatible}.
\item \textbf{Extract:} Compute architecture-aware branch products $\{M_\ell^A\}$, $\{M_\ell^B\}$.
\item \textbf{Fingerprint:} Compute centered residual signatures $\{\phi_\ell^A\}$, $\{\phi_\ell^B\}$.
\item \textbf{Align:} Apply Hungarian matching to recover block correspondences.
\item \textbf{Aggregate:} Compute lineage score $\mathcal{L}(A,B)$.
\item \textbf{Calibrate:} Compute z-score $Z(A,B)$.
\item \textbf{Decide:} Return verdict:
\end{enumerate}
\begin{equation}
\mathrm{Verdict} = \begin{cases}
\textsc{Descendant} & \text{if } Z(A,B) > \tau_{\mathrm{upper}} \\
\textsc{Non-Descendant} & \text{if } Z(A,B) < \tau_{\mathrm{lower}} \\
\textsc{Inconclusive} & \text{otherwise}
\end{cases}
\end{equation}

\subsection{Failure Modes}

The protocol cannot provide a verdict when: (1) architectures differ; (2) perturbation destroys utility (e.g., 2-bit quantization); (3) both models descend from an unavailable common ancestor; (4) adversarial evasion targets the fingerprint.

\subsection{Lineage Chain and Branching Tree}

To test multi-generational ancestry recovery, we construct a branching tree: root $C_1$, siblings $C_2, C_3$ fine-tuned from $C_1$, grandchildren $C_4 = C_2 \to \mathrm{FT}$, $C_5 = C_3 \to \mathrm{FT}$. The lineage score recovers correct ancestry: for $C_4$, ranking by $\mathcal{L}(C_4, \cdot)$ is parent ($0.96$) $>$ grandparent ($0.91$) $>$ uncle ($0.87$) $>$ cousin ($0.85$) $>$ independent ($0.08$). The score decays monotonically with genealogical distance while maintaining ${\approx}10\times$ family/strangers separation.

\section{Per-Architecture Details}\label{app:architectures}

\subsection{GPT-2 Scaling}

Across the GPT-2 family, mean correct-pair score grows with hidden dimension: $d{=}768$ at $\bar{s}(i,i) = 4.18$, $d{=}1024$ at $5.46$, $d{=}1280$ at $6.98$, $d{=}1600$ at $7.77$. The scaling matches the $\Theta(\sqrt{d})$ prediction within $\le 8\%$.

\subsection{Alternative Methods Comparison}

Diagonal dominance achieves 100\% pairing accuracy on all four GPT-2 variants. Frobenius norm matching degrades from 67\% (GPT-2) to 47\% (GPT-2-large). Singular-value distance performs at chance (0--8\%). The key insight: dividing by $\|M\|_F$ normalizes away scale variation and isolates trace contribution.

\subsection{Per-Kind Lineage Baseline Breakdown}\label{app:per-kind-baselines}

Table~\ref{tab:per-kind-baselines} reports the mean lineage score of each baseline broken down by descendant kind (fine-tune, fine-tune to a new target, noise, prune, quantize) and non-descendant kind (different-seed same-task, distilled). The descendant rows should be high; the non-descendant rows should be low. The aggregate AUROC reported in Table~\ref{tab:baselines} hides two structural failures that this breakdown exposes:

\begin{itemize}[nosep]
\item \textbf{CKA and SVCCA fail on distilled non-descendants.} CKA scores distilled students at $0.83$ and SVCCA at $0.89$---only $1.2\times$ and $1.13\times$ below the corresponding fine-tune descendants ($0.999$ and $0.9997$). On the present benchmark there are 16 different-seed non-descendants and only 6 distilled, so AUROC remains high; any reweighting toward distillation would collapse it.
\item \textbf{IPGuard scores noise-perturbed descendants at $\mathbf{0.000}$.} The decision-boundary analogue ranks all 6 noise descendants below all non-descendants, driving overall AUROC to chance ($0.500$).
\end{itemize}

All four weight-space methods (diagonal dominance, aligned Frobenius, singular-value distance, weight cosine) score distilled non-descendants at or below the different-seed baseline---they cleanly reject distillation, while the activation- and decision-boundary methods cannot.

\begin{table*}[t]
\centering
\small
\begin{tabular}{@{}lccccccc@{}}
\toprule
 & \multicolumn{5}{c}{Descendants (should be high)} & \multicolumn{2}{c}{Non-descendants (should be low)} \\
\cmidrule(lr){2-6} \cmidrule(lr){7-8}
Method & FT & FT-new tgt & Noise & Prune & Quant & Diff-seed & Distilled \\
\midrule
\textbf{Diagonal Dominance (ours)} & 0.391 & 0.326 & 0.400 & 0.359 & 0.404 & 0.029 & \textbf{0.030} \\
Aligned Frobenius & $-0.036$ & $-0.301$ & $-0.026$ & $-0.204$ & $-0.007$ & $-1.333$ & $-1.300$ \\
Singular Value Dist.\ & $-0.004$ & $-0.023$ & $-0.003$ & $-0.018$ & $-0.001$ & $-0.030$ & $-0.032$ \\
Weight Cosine & 0.999 & 0.956 & 1.000 & 0.979 & 1.000 & 0.096 & 0.109 \\
SVCCA & 1.000 & 0.953 & 1.000 & 0.993 & 1.000 & 0.881 & \emph{0.885} \\
CKA & 0.999 & 0.760 & 1.000 & 0.986 & 1.000 & 0.820 & \emph{0.827} \\
IPGuard (regr.) & 0.888 & 0.110 & \emph{0.000} & 0.626 & 0.993 & 0.632 & 0.642 \\
\bottomrule
\end{tabular}
\caption{Per-kind mean lineage scores on the 52-pair MLP benchmark. Each method uses its native scale; what matters is the relative gap between the rightmost columns (non-descendants) and the leftmost five (descendants). Italicized entries highlight failure modes: SVCCA and CKA give distilled non-descendants near-descendant scores; IPGuard scores noise-perturbed descendants at exactly $0.000$. All four weight-space methods (including ours) keep distilled below or at the diff-seed baseline. $n{=}6$ per descendant kind, $n{=}16$ diff-seed, $n{=}6$ distilled.}
\label{tab:per-kind-baselines}
\end{table*}

\subsection{ResNet Factorization}

For Bottleneck blocks, the naive $W_3 W_1$ achieves chance-level accuracy. The triple product $W_3 W_2 W_1$ recovers:
\begin{itemize}[nosep]
\item ResNet-50/layer3 ($n{=}5$): 100\%, AUC 1.000
\item ResNet-101/layer3 ($n{=}22$): 100\%, AUC 0.995
\item ResNet-152/layer3 ($n{=}35$): 91\%, AUC 0.964
\end{itemize}

\subsection{CIFAR-10 ResNet-18 Benchmark}

We train two CIFAR-10 ResNet-18 references and three independents. Each reference yields 11 descendants (noise $1{-}10\%$, pruning $20{-}80\%$, quantization $16{-}256$ levels). Results: AUROC$=$1.000, TPR@1\%FPR$=$100\%. Descendant scores: $\mathcal{L} \in [0.695, 1.000]$; independent baseline: $\mathcal{L}_{\max} = 0.004$.

\subsection{Branching Ancestry Recovery}

On a synthetic tree $C_1 \!\to\! \{C_2, C_3\}$, $C_2 \!\to\! C_4$, $C_3 \!\to\! C_5$, the lineage score ranks candidates correctly:
\begin{align*}
\mathcal{L}(C_4, C_2) &= 0.961 \;>\; \mathcal{L}(C_4, C_1) = 0.910 \\
&>\; \mathcal{L}(C_4, C_3) = 0.867 \;>\; \mathcal{L}(C_4, C_5) = 0.850 \\
&\gg\; \mathcal{L}(C_4, I_k) \approx 0.08.
\end{align*}
The score decays monotonically with genealogical distance while maintaining ${\approx}10\times$ family/strangers separation.

\section{Attack and Robustness Details}\label{app:attacks}

\subsection{Gradient-Based Suppression Attack}

An adversary optimizes perturbation $\Delta$ to minimize lineage score subject to utility:
\begin{equation}
\min_{\Delta} \;\mathcal{L}_{\mathrm{cos}}(A, A + \Delta) \;+\; \lambda \cdot \|f_{A+\Delta}(X) - f_A(X)\|_2^2
\end{equation}

\begin{table}[h]
\centering
\small
\begin{tabular}{@{}rccc@{}}
\toprule
$\lambda$ & Final $\mathcal{L}$ & Eval loss & Verdict \\
\midrule
$0$         & 0.053 & 39.5 & utility destroyed \\
$10^{-2}$   & 0.054 & 0.77 & +12\% loss \\
$10^{-1}$   & 0.083 & 0.70 & +1.5\% loss \\
$\ge 1$     & 0.91+ & 0.69 & utility preserved \\
\bottomrule
\end{tabular}
\caption{Pareto frontier: suppressing $\mathcal{L}$ to chance level requires $\ge$12\% utility loss.}
\end{table}

\subsection{Reparameterization Invariances}

\textbf{Permutations:} $W_{\mathrm{in}} \leftarrow P W_{\mathrm{in}}$, $W_{\mathrm{out}} \leftarrow W_{\mathrm{out}} P^\top$ leaves $M = W_{\mathrm{out}} W_{\mathrm{in}}$ unchanged.

\textbf{Rescaling:} $W_{\mathrm{in}} \leftarrow c W_{\mathrm{in}}$, $W_{\mathrm{out}} \leftarrow W_{\mathrm{out}}/c$ preserves $M$ exactly.

\textbf{Orthogonal rotations:} $R^\top R = I$ preserves $M$ but breaks model function (ReLU non-commutative).

All three yield $\mathcal{L} = 1.0$ to numerical precision---these are invariances, not evasion paths.

\section{Practical Applications}\label{app:applications}

The fingerprint enables applications beyond forensic verification.

\subsection{Training Quality Assurance}

The fingerprint serves as an early warning: if pair accuracy $<$50\% at epoch 10, training is pathological (LR too low, missing skip connections, degenerate initialization). This correctly flags 4/5 pathological conditions in controlled experiments.

\subsection{Zero-Knowledge Ownership Proofs}

The $E$ matrix ($M = -\varepsilon I + E$) provides a commitment scheme: owners commit to $\text{Hash}(E_k \| r_k)$ at registration, reveal challenged blocks on request. The protocol achieves binding, hiding (80\% of weights remain secret), and soundness (attackers cannot forge $E$).

\subsection{Compression Auditing}

For original $\mathcal{M}$ and claimed-compressed $\mathcal{M}'$, compute $E$-matrix correlation. Quantization ($r > 0.97$), pruning ($r > 0.75$ at 70\%), and fine-tuning ($r > 0.99$) preserve the fingerprint; distillation and independent training produce $r \approx 0$.
